\section{结论}

本文从数据、部署和评测三个层面系统推进了实用 GUI 智能体的发展。我们提出 CSRS 与自进化训练流水线，以更低成本构建高质量 GUI 轨迹与知识数据；基于此训练的 Step-GUI 系列模型在多个 GUI 基准上取得了领先或有竞争力的表现，并保留了较强的通用多模态能力。

在部署层面，我们提出 GUI-MCP，为 LLM 与异构设备之间的交互提供了统一、分层且注重隐私的协议；在评测层面，我们提出 AndroidDaily，把 GUI 智能体的评价进一步拉近到真实移动日常使用场景。整体而言，这些工作缩小了 GUI 智能体“研究原型”和“可靠日常助手”之间的距离。
